\begin{titlepage}
\phantomsection
\label{sec:oc133-title-page}
\thispagestyle{empty}
\definecolor{ocTitleBlue}{HTML}{173B57}
\definecolor{ocTitleGold}{HTML}{A77D2A}
\begin{tikzpicture}[remember picture,overlay]
\draw[ocTitleBlue,line width=0.75pt]
  ([xshift=1.35cm,yshift=-1.35cm]current page.north west)
  rectangle
  ([xshift=-1.35cm,yshift=1.35cm]current page.south east);
\draw[ocTitleGold,line width=0.35pt]
  ([xshift=1.55cm,yshift=-1.55cm]current page.north west)
  rectangle
  ([xshift=-1.55cm,yshift=1.55cm]current page.south east);
\end{tikzpicture}
\begin{center}
\vspace*{1.0cm}
{\sffamily\Large\color{ocTitleBlue}\textsc{Ontology of Continua}}\\[0.38cm]
{\sffamily\Huge\bfseries Core~1.3.3}\\[0.28cm]
{\sffamily\Large Master Monograph}\\[0.75cm]
{\color{ocTitleGold}\rule{0.42\textwidth}{0.5pt}}\\[0.95cm]

{\Large Alexander Yashin}\\[0.2cm]
{\normalsize Independent Researcher, Leipzig/Halle, Germany}\\[0.2cm]
{\normalsize \href{https://orcid.org/0009-0008-6166-0914}{ORCID 0009-0008-6166-0914}}\\[1.0cm]

{\normalsize Version v1.3.3; release date: 1 May 2026}\\[0.35cm]
{\normalsize Version DOI: \href{https://doi.org/10.5281/zenodo.17899134}{10.5281/zenodo.17899134}}\\[0.2cm]
{\normalsize Zenodo record: \href{https://doi.org/10.5281/zenodo.17899134}{https://doi.org/10.5281/zenodo.17899134}}\\[0.2cm]
{\normalsize Concept DOI: \href{https://doi.org/10.5281/zenodo.17899134}{10.5281/zenodo.17899134}}\\[0.2cm]
{\normalsize GitHub release: \href{https://github.com/alexanderyashin/ontology-of-continua-core-main/releases/tag/v1.3.3}{v1.3.3}}\\[0.9cm]

\vfill

% DEDICATION_REQUIRED_DO_NOT_REMOVE
% OC_CORE_PUBLICATION_DEDICATION
{\small\itshape Dedicated to my dear wife Maria, without whom this work would have been impossible.}\\[0.8cm]

\begin{minipage}{0.86\textwidth}
\centering\scriptsize
This public monograph is the canonical OC Core~1.3.3 scientific artifact.
Earlier source witnesses and archived releases are provenance history only; they are not the current release identity.
\end{minipage}
\end{center}
\end{titlepage}

\clearpage
\noindent\textbf{Acknowledgements.}
The author thanks the reviewers and critics whose questions, objections, and
suggestions contributed to the development of the Ontology of Continua and to
the refinement of this release. Their criticism helped sharpen the claim
boundaries, improve the reader route, expose weak presentation choices, and
force clearer separation between model claims, proof routes, numerical evidence,
prior-art comparison, and publication form. Acknowledgement records review
pressure and intellectual contribution to the development of the work; it does
not imply authorship, endorsement, publication approval, responsibility for the
theory, or agreement with any claim promoted in the manuscript.
\begin{itemize}
    \item G.~V.~Apostolov.
    \item Eduard Fadeev.
    \item Gennady Alekseevich Nosov.
    \item Sergey Shpadyrev.
    \item Stanislav Tsukrov.
\end{itemize}

\clearpage
\begin{abstract}
Ontology of Continua (OC) Core~1.3.3 is the bounded external-review release of
the OC model core. It presents a typed account of continuants, realizations,
liveness, death, residue, rebirth, morphisms, generalized boundaries, hybrid
operators, cycle modes, historical and effective dimension, and K-level
witnesses under explicit assumptions. The scientific task of this manuscript is
not to replace every domain science with a slogan. Its task is to give reviewers
a precise object whose definitions, proof routes, finite semantic checks,
figures, tables, numerical anchors, comparator boundaries, and failure
conditions can be inspected.

The monograph is written as one continuous scientific manuscript rather than as
a prior volume followed by an update packet. The 1.3.3 additions are integrated
into the model, proof, evidence, comparison, limitation, and reviewer-boundary
chapters. The reader should therefore treat the text as the canonical public
spine of the release: it teaches the model first, then binds theorem routes,
Lean subset evidence, finite-model witnesses, bounded target-blind replay QA,
prior-art positioning, and adversarial-review material to the claim surface.

The current maturity of OC Core is deliberately stated as bounded. The release
is mature enough to be read as a coherent scientific manuscript and to be
checked against formal and executable artifacts. It is not a declaration of
complete scientific coverage, complete numerical closure, or unbounded
comparative victory over contemporary science. Where a theorem route, proof sheet,
Lean declaration, finite semantic case, replay row, comparator row, or
negative-control route supports a claim, the claim is promoted. Where that
support is incomplete, the text records a limitation or future obligation
instead of widening the public wording.

The manuscript contains the full monograph body, a compact article route, a
methods and reproducibility route, reviewer-response material, proof and
formalization sections, finite-model and numerical evidence anchors, figure and
table material, source-trace appendices, and practical-use boundaries. Figures
and tables are retained as scientific material; formulas and theorem statements
are retained in the mathematical spine; machine-readable registers remain in
the evidence package where they can be audited without becoming the narrative
structure of the manuscript.

The principal limitation is part of the method. A reader should evaluate whether
the typed model, proof governance, finite semantics, bounded replay evidence,
prior-art comparison, and adversarial-response route form a coherent model-core
contribution under stated assumptions. Future releases must add evidence,
mechanization, comparator work, or domain evidence-boundary review before they widen the
public claim surface. Version~1.3.3 therefore offers a publication-grade basis
for external review, not a claim that every future scientific obligation has
already been discharged.
\end{abstract}

\clearpage
\section*{Version 1.3.3 Release Delta}
\addcontentsline{toc}{section}{Version 1.3.3 Release Delta}
Version~1.3.3 is the release in which the OC Core corpus is reorganized from
scattered source witnesses and process records into a reviewable scientific
package. The visible delta is not merely a new archive number: the release
strengthens the typed foundation, makes the claim boundary explicit, integrates
the theorem route with public proof sheets, and separates reader-facing prose
from machine evidence.

The model delta includes clearer treatment of carriers, realizations, liveness,
death, residue, rebirth, morphisms, generalized boundaries, hybrid operators,
cycle modes, historical and effective dimension, and K-level witnesses. These
additions are presented as a bounded model-core grammar rather than as an
unrestricted theory of every phenomenon.

The verification delta adds and consolidates a Lean-checked subset, structured
proof sheets, finite semantic checks, finite witness interpretation, bounded
target-blind replay QA, comparator and prior-art positioning, negative-control
and falsifier language, and adversarial-review response material. Where a
stronger claim is not yet earned, the release records the limitation instead of
hiding it inside a source-control record.

The editorial delta is equally important. Figures, tables, formulas, numeric
anchors, appendices, and evidence routes are kept in the publication corpus;
machine coverage maps remain coverage and trace data only. The public
documents are therefore built from the human manuscript hierarchy and curated
payload sources, while source bindings and machine evidence indexes stay in the
review manifest.

\clearpage
\section*{Reader Contract}
\addcontentsline{toc}{section}{Reader Contract}
Dear reader, this monograph is the complete public manuscript for OC
Core~1.3.3. It is not an internal routing memo, not a build log, and not a
one-to-one printout of evidence-node records. It exists so that a careful
reader can understand the model before opening the machine-readable evidence
package.

The work may be useful to several groups: scientists and philosophers testing
the continuum model, formal-methods readers checking proof obligations, systems
theorists comparing cross-domain structure, reproducibility reviewers replaying
finite and numerical evidence, journal editors looking for a compact claim
boundary, and hostile reviewers searching for overclaim, equivalence to prior
art, or missing falsifiers.

The recommended route is simple. Read the abstract and release delta first.
Use the table of contents to enter the main model chapters. The conceptual
model and formulas live in the scientific body; the proof route and Lean subset
live in theorem, proof, and formalization sections; numeric evidence and
target-blind replay QA live in the methods and evidence route; figures and
tables live in the monograph and its figure and evidence appendices; prior-art
comparison and reviewer objections live in the comparison and attack-response
sections.

Read the limitations as part of the claim, not as a footnote. A statement is
promoted only where its assumptions, proof or evidence class, comparator
boundary, and reopening condition are visible. Claims about complete scientific
coverage, complete numerical closure, or unbounded comparative victory over
contemporary science are not promoted by this release.

\clearpage
\noindent\textbf{Keywords.}
Ontology of Continua; typed model core; formal methods; proof governance;
finite semantic checks; target-blind replay QA; reproducible research;
scientific release engineering.

\clearpage
\noindent\textbf{Reproducibility and citation note.}
\begin{itemize}[leftmargin=1.6em,nosep]
\item Current version DOI: \href{https://doi.org/10.5281/zenodo.17899134}{10.5281/zenodo.17899134}.
\item Current Zenodo record: \href{https://doi.org/10.5281/zenodo.17899134}{https://doi.org/10.5281/zenodo.17899134}.
\item GitHub release tag: \href{https://github.com/alexanderyashin/ontology-of-continua-core-main/releases/tag/v1.3.3}{v1.3.3}.
\end{itemize}

\clearpage
